\section{Windows PowerShell}
\label{sec:powershell}

\subsection{Background}

Most of the project is operated by typing commands rather than by clicking. A program that
accepts typed commands is called a \term{terminal} or \term{command line}; the one that every
Windows computer already contains is \term{Windows PowerShell}. A command is one line of text.
When you press \key{Enter}, PowerShell carries the line out, prints whatever the command
reports, and then waits for the next line. Nothing in this guide requires PowerShell to be
started ``as administrator'', and you should not start it that way.

\subsection{Open a PowerShell window}

\step{Open PowerShell}\label{st:openps}
Click the \ui{Start} button (or press the \key{Windows} key), type \code|powershell|, and wait
for the search results. Click the result named \ui{Windows PowerShell} whose description is
\emph{App}. Do \emph{not} choose \emph{Windows PowerShell ISE}, \emph{Windows PowerShell (x86)}
or \emph{Command Prompt}; they behave differently.

A window with a dark background opens. On Windows~11 it is usually a window of the
\emph{Terminal} app with a tab labelled \emph{Windows PowerShell}; on Windows~10 it is a blue or
black window titled \emph{Windows PowerShell}. Both are the same program.

\begin{out}[What the new window shows]
Windows PowerShell
Copyright (C) Microsoft Corporation. All rights reserved.

PS C:\Users\ttrojan>
\end{out}

\checkpoint{The last line begins with \texttt{PS}, shows a folder name, and ends with
\texttt{>} followed by a blinking cursor. On your computer the folder is
\texttt{C:\bs Users\bs} followed by your own Windows user name.}

\subsection{Read the prompt}

The last line of the window is the \term{prompt} (Figure~\ref{fig:prompt}). It has three parts.
\texttt{PS} says that PowerShell is ready for a command. The path after it, here
\code|C:\Users\ttrojan|, is the \term{current folder}: commands that mention a file look for it
in this folder unless told otherwise. The sign \code|>| marks the place where your typing
appears. The prompt is printed by PowerShell; you never type it, and it is not part of any
command in this guide.

\begin{figure}[H]
\centering
\begin{tikzpicture}[font=\ttfamily\small, seg/.style={fill=black!88, text=white, inner xsep=0pt, inner ysep=7pt, outer sep=0pt, anchor=west},
  lab/.style={font=\rmfamily\footnotesize, align=center, text=black}]
\node[seg, inner xsep=0pt] (pad) at (0,0) {\hspace{7pt}};
\node[seg, text=green!70!white] (a) at (pad.east) {(.venv)\vphantom{Ag>}};
\node[seg] (b) at (a.east) {\ PS\ \vphantom{Ag>}};
\node[seg] (c) at (b.east) {C:\textbackslash ee105\textbackslash camera-to-latex>\vphantom{Ag}};
\node[seg] (d) at (c.east) {\ python -m pytest -q\hspace{7pt}\vphantom{Ag>}};
\draw[decorate, decoration={brace, mirror, amplitude=4pt}] ($(a.south west)+(0,-0.08)$) -- ($(a.south east)+(0,-0.08)$);
\draw[decorate, decoration={brace, mirror, amplitude=4pt}] ($(b.south west)+(0.12,-0.08)$) -- ($(b.south east)+(-0.12,-0.08)$);
\draw[decorate, decoration={brace, mirror, amplitude=4pt}] ($(c.south west)+(0,-0.08)$) -- ($(c.south east)+(0,-0.08)$);
\draw[decorate, decoration={brace, mirror, amplitude=4pt}] ($(d.south west)+(0.15,-0.08)$) -- ($(d.south east)+(-0.25,-0.08)$);
\node[lab, anchor=north] at ($(a.south)+(-0.35,-0.95)$) {active environment\\(Section~\ref{sec:python})};
\draw[black!50] ($(a.south)+(0,-0.3)$) -- ($(a.south)+(-0.35,-0.95)$);
\node[lab, anchor=north] at ($(b.south)+(0,-0.3)$) {ready};
\node[lab, anchor=north] at ($(c.south)+(0,-0.3)$) {current folder};
\node[lab, anchor=north] at ($(d.south)+(0,-0.3)$) {the command you typed};
\end{tikzpicture}
\caption{The parts of a PowerShell line as it will look later in this guide. Everything to the
left of the command is the prompt, which PowerShell prints by itself.}
\label{fig:prompt}
\end{figure}

\subsection{Type, paste and run commands}

\step{Run a first command}
Click once inside the window, type the following line, and press \key{Enter}.
\begin{ps}
Get-Location
\end{ps}
\begin{out}
Path
----
C:\Users\ttrojan

PS C:\Users\ttrojan>
\end{out}
\checkpoint{PowerShell printed the current folder and then a new prompt. The new prompt is the
sign that the command has finished; a command is still running for as long as no prompt has
appeared.}

\step{Paste a command instead of typing it}
Typing errors are the commonest cause of failed commands, so pasting is preferable. In this PDF,
select the line in the blue box below by dragging the mouse from its first character to its
last (the title bar of the box must not be selected). Press \key{Ctrl}+\key{C} to copy. Click in
the PowerShell window and press \key{Ctrl}+\key{V}, or click the right mouse button, to paste.
Compare the pasted line with the box, then press \key{Enter}.
\begin{ps}
Get-ChildItem
\end{ps}
\checkpoint{PowerShell lists the folders in your user folder (\file{Desktop}, \file{Documents},
\file{Downloads}, and so on) and prints a new prompt.}

\begin{iffail}[If a warning about ``multiple lines'' appears when you paste]
You copied more than one line. Choose \ui{Cancel}, copy the single line again, and paste it. A
pasted line break acts like the \key{Enter} key, so a paste of two lines runs the first command
before you have seen the second.
\end{iffail}

\step{Clear the window and recall a command}
Type \code|cls| and press \key{Enter}: the window is emptied (nothing is undone; only the
display is cleared). Now press the \key{$\uparrow$} key once: PowerShell brings back the
previous command so that it can be run again or corrected. Press \key{Esc} to discard the
recalled line.
\checkpoint{The window shows a single prompt and an empty line after it.}

\subsection{Rules of conduct at the command line}

Table~\ref{tab:psrules} states the habits on which the rest of the guide relies.

\begin{table}[H]
\centering
\begin{tabular}{L{4.5cm}L{7.1cm}}
\toprule
\multicolumn{1}{c}{\hd{Rule}} & \multicolumn{1}{c}{\hd{Reason}} \\
\midrule
Enter one command at a time, and wait for the prompt before the next &
  A later command usually depends on the result of an earlier one \\
Read the output before continuing &
  Text in red is an error message; the command did not do its work \\
Stop a running command with \key{Ctrl}+\key{C}, not by closing the window &
  The program is told to stop and can tidy up \\
Do not click inside a classic PowerShell window while a command runs &
  A click starts \emph{selection}, which freezes the output until \key{Esc} is pressed; the word
  \emph{Select} appears in the title bar \\
Open a new window when the guide says so &
  Newly installed programs are found only by windows opened after the installation \\
Close a window with the command \code|exit| or the \ui{X} button &
  Either is safe once a prompt is showing \\
\bottomrule
\end{tabular}
\caption{Working habits for PowerShell.}
\label{tab:psrules}
\end{table}

Whenever a later section asks for ``a new PowerShell window'', it means: carry out
Step~\ref{st:openps} again. From Section~\ref{sec:python} onward, a new window also needs two
commands that return it to the project; they are written out in full every time they are
needed, and they are collected in Section~\ref{sec:opening}.
